% answers_short.tex —— 2009 年·普通高等学校招生考试·重庆卷·文科 · 简短答案速查
% 由 tools/make_short_answers.py 自 answers.tex 的 \gksolhead/\gkanswer 提取，
% 并与 2009/2009_数学真题_重庆文_答案要点.md 逐题比对后生成，勿手改（改源 MD 后重跑）。

\begin{gknotice}
\textbf{使用说明}：本节只给最终答案便于快速核对，完整推导见解析版。答案为编写者独立推导，\textbf{非官方标准答案}。
\end{gknotice}

\gkpart{选择题}
\begin{center}\begin{tabular}{ccccccccccc}
\toprule
题号 & 1 & 2 & 3 & 4 & 5 & 6 & 7 & 8 & 9 & 10 \\
\midrule
答案 & A & B & D & D & A & C & C & B & C & B \\
\bottomrule
\end{tabular}\end{center}

\gkpart{填空题}
\begin{itemize}[leftmargin=2.6em]
  \item[\textbf{11.}] $\{2,4,8\}$。
  \item[\textbf{12.}] $2$。
  \item[\textbf{13.}] $72$。
  \item[\textbf{14.}] $\sqrt5$。
  \item[\textbf{15.}] $[\sqrt2-1,1)$。
\end{itemize}

\gkpart{解答题}
\begin{itemize}[leftmargin=2.6em]
  \item[\textbf{16.}] （1）$\omega=\dfrac32$，最小正周期 $\dfrac{2\pi}3$；（2）$\left[\dfrac\pi4+\dfrac{2k\pi}3,\dfrac{7\pi}{12}+\dfrac{2k\pi}3\right]$，$k\in\mathbf Z$。
  \item[\textbf{17.}] （1）$\dfrac{899}{900}$；（2）$\dfrac4{45}$。
  \item[\textbf{18.}] （1）$\dfrac{2\sqrt5}5$；（2）$\sqrt2$。
  \item[\textbf{19.}] （1）$a\leqslant-\sqrt3$ 或 $a\geqslant\sqrt3$；（2）增区间 $(-\infty,-1)$，$\left(-\dfrac13,+\infty\right)$，减区间 $\left(-1,-\dfrac13\right)$。
  \item[\textbf{20.}] （1）$x^2-\dfrac{y^2}4=1$；（2）$\sqrt{10}+1$，$M\left(\dfrac{4\sqrt2-\sqrt5}3,\dfrac{4(\sqrt5-\sqrt2)}3\right)$。
  \item[\textbf{21.}] （1）$b_1=4$，$b_2=\dfrac{17}4$，$b_3=\dfrac{72}{17}$；（2）证明见解析；（3）证明见解析。
\end{itemize}
